\documentclass[11pt,a4paper]{letter}
\usepackage[utf8]{inputenc}
\usepackage[T1]{fontenc}
\usepackage{mlmodern}
\usepackage{amsmath}
\usepackage{amsfonts}
\usepackage{amssymb}
\usepackage{graphicx}
\usepackage{geometry}
\usepackage{csquotes}
\usepackage[final]{microtype}
\geometry{margin=1in}

\address{Aniara Lyn Lauron \\ Gåseholmvej 61 \\ Herlev, Denmark \\ Email: aniaralyn.lauron@gmail.com \\ Phone: +45 91 80 63 94}
\date{March 7, 2026}

\begin{document}

\begin{letter}{Admissions Committee \\ Absalon University College \\ Campus Kalundborg 4, 4400 Kalundborg\\ Denmark}

\opening{Dear Admissions Committee,}

I am writing to express my strong interest in the Bachelor of Engineering in Biotechnology programme at Absalon University College, starting in the fall of 2026. With a prior Bachelor of Science in Pharmacy from the University of Southern Philippines, I am eager to build on my scientific foundation and gain the engineering expertise needed to develop biotechnological solutions for health, food and the environment.

My pharmacy training provided me with a deep understanding of chemistry, biology and experimental methods. For my bachelor thesis, \enquote{The Effect of Hagonoy Leaves Extract on the Inflamed Paws of Male Albino Rabbits}, I performed chemical extractions, managed animal experiments, and analysed results statistically. I also carried out toxicology studies, calculating $LD_{50}$ values and monitoring drug exposure are experiences that sharpened my attention to detail and instilled an appreciation for ethical laboratory practice.

Outside the classroom I have held roles that strengthened my teamwork and communication skills as a pharmacy technician at Optum Global Solutions Inc. and Collabera Technologies and freelance English instructor and most recently, hotel assistant at The Ellen Group. Each position taught me to problem solve, collaborate across disciplines and adapt quickly, qualities that are essential in interdisciplinary engineering projects. I am particularly genuine about helping others and motivated to do better in any role I take on, which I believe will serve me well in the collaborative environment of Absalon University College. I also build strong relationships with colleagues and customers, which I hope to continue in the biotech community. I like to engage and make people feel satisfied with the work I do, and I am excited to bring that enthusiasm to the field of biotechnology.

Living and working in the Philippines, South Korea, the Netherlands and Denmark has given me a global perspective and exposed me to many engineers. Their enthusiasm for technology first sparked my curiosity about how biological systems can be harnessed for industrial purposes. This motivated me to undertake hands on technical projects on my own time: in 2024 I designed and built an Arduino driven automated plant irrigation system, writing firmware in C++ and integrating sensors and actuators. In 2026 I am tackling a custom piano computer workstation, exercising skills in woodworking, design and assembly. I also made a silver ring in 2023, learning metalworking and engraving along the way.

The Biotechnology programme’s combination of biology, chemistry and science of engineering, together with its focus on using cells as “small factories,” aligns perfectly with my interests. I am particularly drawn to the emphasis on engineering principles for processing materials by microorganisms or plant and animal cells, and I appreciate that the curriculum bridges theory and practice from the very beginning through close cooperation with local companies in Kalundborg. During my pharmacy studies I conducted simple cell growth experiments, but the opportunity at Absalon University College to model such systems mathematically and scale them for production excites me. I am also keen to learn the techniques used in the pharmaceutical and food industries, and to apply them to environmental challenges.

My long term goal is to merge my background in pharmacy with engineering skills to develop safer and more efficient bioproduction processes. After completing the BEng, I plan to find a job degree ideally in bioengineering or bioprocess technology and to contribute to Denmark’s leading biotech companies such as Novo Nordisk or Coloplast. I envision working in research and development on drug manufacturing, bringing life improving products to market. I am excited about the chance to contribute my unique perspective and to grow within Denmark’s vibrant biotechnology community.

Thank you for considering my application. I look forward to the possibility of discussing how my experiences and aspirations align with your institution’s mission.

\closing{Sincerely, \\
Aniara Lyn Lauron}

\end{letter}

\end{document}